% Kalman Filter Hedge Ratios — research-note PDF.
%
% PROSE AND CODE ARE VERBATIM from the website article:
%   site/src/content/articles/kalman-filter-hedge-ratios.tsx
% The article interpolates its computed results into the prose; every
% interpolated value here is the one the site renders, read from
% data/kalman-filter-hedge-ratios.ts. Code cells are inline in the article.
%
% CHARTS generated by figures/make_kalman_filter_hedge_ratios.py from that
% module.
\documentclass{pyportfolios-note}

\noteshorttitle{Kalman Filter Hedge Ratios}

\begin{document}
\thispagestyle{notecover}

\notemasthead{Quantitative Insights}{Tutorial}{Q3 2026}{}

\notetitle{Kalman Filter\\Hedge Ratios}

\notedeck{Every pairs trade hides a regression inside it. Estimate it once and
you have assumed the relationship never moves.}

\notelede{Every pairs trade hides a regression inside it: the hedge ratio that
turns two prices into one spread. Estimate that ratio once and you have quietly
assumed the relationship never moves — over fifteen years of EWA and EWC it moves
from 0.71 to 1.49. We build a Kalman filter from scratch in NumPy that treats the
hedge ratio as a moving target, then run an honest experiment: identical trading
rules on the Kalman spread and the static spread, with costs, and let the data
pick the winner.}

\notecard
  {Algorithmic Trading}
  {NumPy · Pandas · statsmodels}
  {EWA + EWC (iShares Australia / Canada)}
  {Jan 2010 – Dec 2024}
  {Keeping a pairs hedge ratio honest as the relationship drifts}

\begin{multicols}{2}
\begin{notepipeline}
\item Load EWA \& EWC adjusted closes, 2010–2024 (the classic Chan pair)
\item Fit the two OLS baselines: full-sample beta and rolling 252d beta
      (statsmodels)
\item Build the Kalman filter in NumPy — state [β, α], random-walk transition
\item Compare the three beta paths on one chart
\item Trade the spread: |z| > 2 entry, z-crosses-0 exit, 10 bp per leg
\item Read the honest scoreboard: gross vs net, Kalman vs static
\end{notepipeline}
\end{multicols}

% ======================================= 1 Why static hedge ratios die ======
\noteneedroom{190bp}
\notesection{1.\notenumsep Why static hedge ratios die}

EWA (Australia) and EWC (Canada) are the canonical cointegration pair — two
commodity-heavy developed markets whose daily returns correlate at
\texttt{0.83}. Regress EWC on EWA over the full sample and you get one hedge
ratio for fifteen years: $\beta = 1.57$, with an ADF p-value of \texttt{0.038} on
the residual spread — cointegrated, by the book (Engle–Granger: regress one price
on the other, then unit-root-test the residual). The problem is the word
\textbf{one}. A hedge ratio is not a constant of nature; it is a snapshot of an
economic relationship — commodity mix, currency betas, index composition — and
every one of those drifted between 2010 and 2024. The standard fix is a rolling
window, and it limps for a reason worth understanding: every observation inside
the window carries equal weight, so year-old data moves today's estimate exactly
as much as yesterday's does — and the day an old observation falls out of the
window, the estimate jerks for no economic reason at all. That artefact has a
name: the \textbf{window cliff}.

\notefigure{figures/kalman-filter-hedge-ratios/fig1-prices.png}
  {EWA \& EWC, normalized to 1.0 at 2010-01-04 — related, not identical}
  {The two series track each other broadly, diverging and reconverging
   repeatedly}
  {Source: Author calculations, from the article's computed results.}

\notecallout{Why practitioners care}{The hedge ratio is not a nuisance parameter
— it \textbf{is} the position. Get it wrong by 20\% and your ``market-neutral''
spread carries a 20\% directional stub of EWA. Static betas fail slowly and
invisibly, which is the most expensive way to fail.}

% =================================== 2 Beta as a state, not a constant ======
\noteneedroom{190bp}
\notesection{2.\notenumsep Beta as a state, not a constant}

\begin{multicols}{2}
The Kalman filter starts from a different premise, and it is worth sitting with
for a moment: the ``true'' hedge ratio is something you can never observe
directly — you only see prices, which are noisy evidence about it. So treat the
regression coefficients as \textbf{unobserved states} drifting through time, and
treat each day's prices as one more noisy measurement of where they are. (A GPS
does exactly this with your position; here the hidden position is β.) Two
equations define the model. The \textbf{state equation} says the hedge ratio and
intercept follow a random walk —
$[\beta_t, \alpha_t] = [\beta_{t-1}, \alpha_{t-1}] + \omega_t$ — tomorrow's
relationship is today's, plus noise. The \textbf{observation equation} says
$\mathrm{EWC}_t = \beta_t\,\mathrm{EWA}_t + \alpha_t + \varepsilon_t$. Two
variances close the model: the state noise
$Q = \frac{\delta}{1-\delta}\,I$ with $\delta = 10^{-5}$ (the standard
parameterization from Chan), and observation noise $R = 10^{-3}$.

Delta is the single real knob, and it replaces the window size entirely: it is a
\textbf{forgetting rate}. Larger δ lets the states wander faster (adaptive but
noisy); smaller δ pins them down (smooth but laggy); δ = 0 collapses to recursive
least squares — a static beta refined forever. Both values here are textbook
defaults, deliberately not tuned on this sample: tune δ to the backtest and you
are optimizing the strategy through the back door.
\end{multicols}

% ================================ 3 The filter in five lines of NumPy =======
\noteneedroom{190bp}
\notesection{3.\notenumsep The filter in five lines of NumPy}

No library, no black box — the whole filter is a predict step and an update step,
looped over the sample. Predict: with a random-walk transition the state estimate
is unchanged and its covariance grows by \texttt{Q} — a day passes, so the filter
becomes a little less sure of where β is. Update: compare the observed EWC to the
prediction, and shift the states toward the error in proportion to the
\textbf{Kalman gain} — the ratio of state uncertainty to total uncertainty. Big
surprise while unsure of yourself: move a lot. Big surprise while confident: blame
measurement noise, barely move.

\begin{notecode}
def kalman_hedge(x, y, delta=1e-5, r_obs=1e-3):
    q = (delta / (1 - delta)) * np.eye(2)     # trans_cov
    state, p_cov = np.zeros(2), np.eye(2)     # [beta, alpha], diffuse start
    betas, alphas = np.zeros(len(x)), np.zeros(len(x))
    for t in range(len(x)):
        h = np.array([x[t], 1.0])             # observation map
        p_cov = p_cov + q                     # predict   (F = I)
        e = y[t] - h @ state                  # innovation
        s = h @ p_cov @ h + r_obs             # innovation variance
        k = p_cov @ h / s                     # Kalman gain
        state = state + k * e                 # update
        p_cov = p_cov - np.outer(k, h @ p_cov)
        betas[t], alphas[t] = state
    return betas, alphas
\end{notecode}

The gain is the elegance. When the filter is uncertain (large \texttt{p\_cov}), it
learns aggressively from each observation; once confident, new data barely moves
it — unless \texttt{Q} keeps injecting doubt, which is exactly what lets β keep
adapting forever. An exponentially-weighted regression, derived from first
principles rather than picked from a menu of window sizes.

% ============================================= 4 Three betas, one pair ======
\noteneedroom{190bp}
\notesection{4.\notenumsep Three betas, one pair}

Slice off the first year (rolling-OLS warm-up, Kalman burn-in) and compare the
three estimators of the same quantity. The static line says the answer is 1.57,
forever. The rolling OLS swings between \texttt{0.19} and \texttt{2.75} — whipping
around every regime change a full window late. The Kalman path stays in a far
narrower band, moving early and smoothly: no cliff, because no window.

\notefigure{figures/kalman-filter-hedge-ratios/fig2-beta.png}
  {The EWC\textasciitilde EWA hedge ratio, three ways — Kalman filter vs rolling 252d OLS vs full-sample OLS}
  {The Kalman β drifts smoothly between roughly 0.7 and 1.5; the rolling OLS β
   oscillates between about 0.2 and 2.7}
  {Source: Author calculations, from the article's computed results.}

Note what the rolling estimator does around 2020–2021: the COVID shock enters the
window, distorts the regression for exactly 252 trading days, then falls out and
the estimate jumps again — two artefacts from one event. The filter digests the
same shock in weeks and moves on.

% ================================================= 5 Trading the spread =====
\noteneedroom{190bp}
\notesection{5.\notenumsep Trading the spread}

Now make the estimator earn its living — a better β is only worth money if it
makes a better spread. The spread is
$\mathrm{EWC}_t - \beta_t\,\mathrm{EWA}_t - \alpha_t$, z-scored on a trailing
60-day window. Rules, identical for both variants: enter long the spread (long
EWC, short β·EWA) when $z < -2$, short when $z > +2$, exit when z crosses zero.
Positions are sized to \$1 gross notional at entry, the Kalman variant re-hedges
the EWA leg to the current β daily, and every unit of traded notional pays 10 bp.
Signals use the close and P\&L starts the next day — no lookahead in the rule. The
static beta itself, of course, is one giant lookahead: it was fit on all fifteen
years, including the future of every trade it takes.

\begin{notecode}
spread = ewc - beta * ewa - alpha              # beta_t, alpha_t from the filter
z = (spread - spread.rolling(60).mean()) / spread.rolling(60).std()

# enter |z| > 2, exit when z crosses 0, hold in between
if p == 0:
    p = 1 if z[t] < -2 else (-1 if z[t] > 2 else 0)
elif p == 1 and z[t] >= 0: p = 0               # long leg reverted
elif p == -1 and z[t] <= 0: p = 0              # short leg reverted
\end{notecode}

\notefigure{figures/kalman-filter-hedge-ratios/fig3-zscore.png}
  {The Kalman-spread z-score with ±2σ entry bands — every excursion is a candidate trade}
  {The z-score oscillates around zero, crossing the entry bands many times across
   the sample}
  {Source: Author calculations, from the article's computed results.}

% ===================================== 6 What the numbers actually say ======
\noteneedroom{190bp}
\notesection{6.\notenumsep What the numbers actually say}

Here is the honest scoreboard, and it is more interesting than a clean win. Gross
of costs, the Kalman spread is the better signal on every risk-adjusted axis —
more Sharpe, half the volatility, half the drawdown. But the filtered spread
mean-reverts \textbf{fast}, so it trades 152 round trips to the static variant's
57 — and at 10 bp per unit of traded notional, that turnover consumes the entire
edge and then some.

\notefigure{figures/kalman-filter-hedge-ratios/fig4-sharpe.png}
  {Sharpe, gross vs net of 10 bp costs — adaptivity wins the signal and loses the implementation}
  {The Kalman variant leads gross (0.55 vs 0.47) but drops below the static
   variant once costs apply (−0.38 vs 0.31)}
  {Source: Author calculations, from the article's computed results.}

\notefigure{figures/kalman-filter-hedge-ratios/fig5-equity.png}
  {Strategy equity, net of 10 bp costs — the cost drag flips the ranking}
  {The static-beta strategy grinds up while the Kalman-beta strategy decays under
   transaction costs}
  {Source: Author calculations, from the article's computed results.}

\begin{notetable}{@{}p{118bp}>{\raggedleft\arraybackslash}p{66bp}>{\raggedleft\arraybackslash}p{50bp}>{\raggedleft\arraybackslash}p{61bp}>{\raggedleft\arraybackslash}p{66bp}>{\raggedleft\arraybackslash}p{61bp}>{\raggedleft\arraybackslash}p{38.8bp}@{}}
\thcell{Variant} & \thcell{Ann ret (net)} & \thcell{Ann vol} & \thcell{Sharpe (net)} & \thcell{Sharpe (gross)} & \thcell{Max DD (net)} & \thcell{Trades} \\
\theadrule
Kalman β (adaptive) & -0.9\% & 2.3\% & -0.38 & 0.55 & -12.6\% & 152 \\ \rowrule
Static β (full-sample) & 1.5\% & 5.0\% & 0.31 & 0.47 & -13.3\% & 57 \\
\end{notetable}

So did the Kalman beta ``improve'' the strategy? As an estimator, unambiguously —
better gross Sharpe, half the drawdown, and no lookahead, against a static
baseline that was handed the answer key. As a net P\&L line at retail costs, no:
the same adaptivity that tracks the relationship also generates signals faster
than 10 bp round trips can pay for. Halve the cost and the gap halves; at
institutional frictions of 1–2 bp the Kalman variant pulls level and ahead.
Estimation quality and implementability are different axes, and a backtest that
reports only one is hiding the other.

\notecallout{Practitioner take}{Use the filter for what it is provably better at:
\textbf{knowing your position}. Desks run Kalman hedge ratios even when entry
signals come from elsewhere, because the alternative is discovering your
directional stub during a drawdown. And before trading the filtered spread, do
the arithmetic this table forces: expected edge per trade must clear cost per
trade — adaptivity raises the trade count, so it raises the bar. δ and R were
fixed at textbook values here; tuning them to fix the net line is how
backtest-overfitting starts. One refinement is legitimate, though: the filter
already computes the spread's fair value \textbf{and its own uncertainty} —
Chan's variant trades the innovation e\textsubscript{t} against its predicted
standard deviation √S\textsubscript{t} instead of a rolling z-score, letting the
model set the entry band and retiring the arbitrary 60-day window.}

\noteneedroom{190bp}
\begin{notereferences}
\item Kalman, R. E. (1960). A New Approach to Linear Filtering and Prediction
      Problems. \textit{Journal of Basic Engineering}, 82(1).
\item Engle, R. F. \& Granger, C. W. J. (1987). Co-integration and Error
      Correction: Representation, Estimation, and Testing.
      \textit{Econometrica}, 55(2), 251–276.
\item Chan, E. (2013). \textit{Algorithmic Trading: Winning Strategies and Their
      Rationale}. Wiley — ch. 3, the EWA/EWC Kalman example and the δ/(1−δ)
      parameterization.
\item Harvey, A. C. (1989). \textit{Forecasting, Structural Time Series Models
      and the Kalman Filter}. Cambridge University Press.
\item Companion notebook: \texttt{kalman-filter-hedge-ratios.ipynb} —
      reproduces every figure from raw data; fully deterministic, no RNG.
\end{notereferences}

\end{document}
